% lemmon-fitch.tex
%
% First draft. Compile with pdflatex. Packages: the AMS bundle, geometry,
% array, and listings -- all part of a standard TeX distribution. The proof
% displays are drawn with a small macro defined below rather than with
% fitch.sty, to keep the source portable; if you would rather use Peter
% Selinger's fitch.sty for the final version, the displays are all in one
% place and easy to convert.
%
% The Haskell excerpts are quoted from the implementation rather than written
% for the paper. Where a definition has been shortened, the omission is marked
% with an ellipsis comment.

\documentclass[11pt]{article}

\usepackage[margin=1.15in]{geometry}
\usepackage{amsmath,amssymb,amsthm}
\usepackage{array}
\usepackage{listings}

\lstset{
  language=Haskell,
  basicstyle=\small\ttfamily,
  keywordstyle=\bfseries,
  commentstyle=\itshape,
  columns=fullflexible,
  showstringspaces=false,
  xleftmargin=1.5em,
  aboveskip=1.2em,
  belowskip=1.2em,
  morekeywords={data,type,where,case,of,let,in,pure}
}

\theoremstyle{plain}
\newtheorem{theorem}{Theorem}
\newtheorem{proposition}[theorem]{Proposition}
\newtheorem{corollary}[theorem]{Corollary}
\newtheorem{lemma}[theorem]{Lemma}
\newtheorem{conjecture}[theorem]{Conjecture}
\theoremstyle{definition}
\newtheorem{definition}[theorem]{Definition}
\newtheorem{example}[theorem]{Example}
\newtheorem{problem}[theorem]{Problem}
\theoremstyle{remark}
\newtheorem{remark}[theorem]{Remark}

% A scope bar for Fitch displays.
\newcommand{\fbar}{\rule[-0.7ex]{0.4pt}{2.6ex}\hspace{0.7em}}
\newcommand{\dep}[1]{\textup{#1}}

\newenvironment{lemmonproof}
  {\begin{center}\begin{tabular}{@{}r@{\hspace{1.5em}}r@{\hspace{0.6em}}l@{\hspace{2.5em}}l@{}}}
  {\end{tabular}\end{center}}

\newenvironment{fitchproof}
  {\begin{center}\begin{tabular}{@{}r@{\hspace{1em}}l@{\hspace{2.5em}}l@{}}}
  {\end{tabular}\end{center}}

\title{Dependency and Scope:\\
       On the Translation Between Lemmon and Fitch Proofs}
\author{Hans Halvorson}
\date{Draft}

\begin{document}
\maketitle

\begin{abstract}
Lemmon-style and Fitch-style natural deduction are usually treated as
notational variants of one another. They are, but the sense in which they are
is more delicate than that description suggests, and the delicacy is
instructive. A Lemmon line carries an explicit and exact record of the
assumptions it rests on; a Fitch line carries no such record, and its
dependencies are bounded, but not determined, by the subproofs enclosing it.
We show that translation from Fitch to Lemmon is total and canonical, and that
the dependency column of a Lemmon proof is therefore redundant. Translation in
the other direction is not \emph{positional}: there are Lemmon proofs whose
lines cannot be carried over one for one, in order, to any Fitch proof. We exhibit two such proofs and
identify the two obstructions responsible. Both, we argue, are consequences of
a single structural difference. Each notation presents a derivation---which is
a tree---as a directed acyclic graph, by writing each line once and citing it
thereafter; but where Lemmon permits a line to be cited by any later line
whatever, Fitch permits citation only along the paths its nesting of subproofs
allows. The obstructions arise exactly when no such nesting accommodates the
sharing the source proof uses, and they dissolve on unfolding, since a graph
without sharing is a tree and a tree constrains nothing. We give a
construction translating any Lemmon proof by way of its unfolding, and then
show that it fails, at exactly one rule. Universal introduction imposes its
side condition on the assumptions a line rests on; Lemmon can state that
exactly and Fitch can state it only about scope, and where the two diverge the
translated proof is incorrect. We exhibit such a proof. The repair is local
and does not move anything: the Lemmon side condition says exactly what is
required to license a substitution, so the offending subderivation may be
renamed to a name occurring nowhere else, whereupon the Fitch condition is
satisfied too. We conjecture that the construction so modified is total, and
identify what a proof of it would need. The construction is implemented, and
its output machine-checked against a regression corpus and sixteen proofs
photographed from student work.
\end{abstract}

\section{Introduction}

Anyone who has taught both notations knows that a Lemmon proof and a Fitch
proof of the same argument look like transcriptions of one another. In the
Lemmon style, each line is annotated with the numbers of the assumptions it
depends on, and a rule such as conditional proof discharges an assumption by
removing its number from the annotation. In the Fitch style, the same
information is carried geometrically: an assumption opens a subproof, drawn as
a vertical scope line, and discharging it closes the subproof. The two
conventions are so evidently doing the same work that the translation between
them is generally left as an exercise, or not mentioned.

It is not quite an exercise. The purpose of this paper is to say exactly what
the translation is, where it is easy, where it is not, and what the difficulty
reveals about the relationship between the two notations.

The central observation is simple enough to state here. In a Lemmon proof, the
dependency set attached to a line is \emph{exact}: it is the set of
assumptions the line actually rests on. In a Fitch proof, the corresponding
information is the set of assumptions of the enclosing subproofs, and this is
merely an \emph{upper bound}. A Fitch line may perfectly well sit inside a
subproof whose assumption it never uses. Every asymmetry we discuss below
follows from this one.

A second observation, less obvious, does most of the work in
Section~\ref{sec:obstructions}. In a Lemmon proof, the textual position of a
line carries no semantic content. Lines may be permuted freely, subject only
to the requirement that a line be written after the lines it cites. In a Fitch
proof, position is semantic: where a line sits determines what may see it, and
a line inside a subproof becomes unavailable once that subproof closes. The
translation from Lemmon to Fitch must therefore invent positional information
that the source does not contain, and it is possible for the source to place
demands that no such invention can satisfy.

A third observation is the paper's main negative result, and we did not
anticipate it. Of the twenty-one rules, exactly one imposes a condition not on
the formulas involved but on the assumptions a line rests on: universal
introduction, whose eigenvariable must be arbitrary. Lemmon can state that
condition exactly. Fitch cannot---it has no record of what a line rests on,
only of what encloses it---and must state the condition about scope instead.
Since scope may properly include the dependency set, the two conditions come
apart, and a correct Lemmon proof may translate to a Fitch proof that breaks
its own rule. Section~\ref{sec:trees} exhibits one.

We work throughout with the rules of \emph{How Logic Works}
\cite{halvorson2020}, a Lemmon-style system with twenty-one rules, and with a
Fitch system chosen to match it rule for rule. Section~\ref{sec:systems} fixes
notation. Section~\ref{sec:fitch-to-lemmon} treats the easy direction and
draws from it the redundancy of the dependency column.
Section~\ref{sec:obstructions} exhibits the two obstructions to a positional
translation in the other direction. Section~\ref{sec:trees} introduces the
graph-theoretic vocabulary needed to say what the obstructions have in common,
locates the difference between the notations in how each constrains the
sharing of lines, gives the construction that translates by unfolding, and
refutes it. Section~\ref{sec:open} states what a proof of the repaired
construction would require, and two questions we have not settled. Section~\ref{sec:machine} reports the machine verification, and
Section~\ref{sec:equivalence} draws the moral.

\section{The two systems}\label{sec:systems}

\begin{definition}
A \emph{Lemmon proof} is a finite sequence of lines. A line is a quadruple
$\langle \Gamma, n, \varphi, j\rangle$ where $n$ is the line's number,
$\varphi$ its formula, $j$ its justification, and $\Gamma$ a set of line
numbers---the \emph{dependency set}. A justification is either
$\mathrm{A}$ (assumption) or a rule together with the numbers of the lines it
cites, all of which must be smaller than $n$.
\end{definition}

The rules fall into two groups. Most of them---modus ponens, conjunction
introduction, universal elimination---simply pool the dependency sets of the
lines they cite. Four of them discharge: conditional proof and reductio each
cite an assumption line together with a line derived from it, and disjunction
and existential elimination each cite one or two such pairs. For these, the
discharged assumption is removed from the pooled set. Thus in

\begin{lemmonproof}
\dep{1}   & (1) & $P$              & A \\
\dep{2}   & (2) & $Q$              & A \\
\dep{1,2} & (3) & $P \wedge Q$     & 1,2 $\wedge$I \\
\dep{1}   & (4) & $Q \to (P \wedge Q)$ & 2,3 CP \\
\end{lemmonproof}

\noindent line (4) discharges assumption (2), which accordingly leaves the
dependency set.

In the implementation a line is a record and a justification an enumeration,
one constructor per rule. The constructors of the discharging rules are worth
noticing: each names an assumption line together with a line derived from it.

\begin{lstlisting}
data ProofLine = ProofLine
  { lineNumber    :: Int
  , formula       :: PredFormula
  , justification :: Justification
  , references    :: Set Int          -- the dependency set
  }

data Justification
  = Assumption
  | MP Int Int                        -- modus ponens, two lines
  | AndIntro Int Int
  | CP Int Int                        -- assumption line, conclusion line
  | RAA Int Int
  | OrElim Int Int Int Int Int        -- disjunction, then two such pairs
  | ExistsElim Int Int Int
  | ForallIntro Int
  -- ... fifteen further rules
\end{lstlisting}

\begin{definition}\label{def:fitch}
A \emph{Fitch proof} is a finite sequence of \emph{items}, where an item is
either a line---a number, a formula and a justification---or a
\emph{subproof}, consisting of an assumption line together with a further
sequence of items. The \emph{scope} of a line is the set of assumptions of the
subproofs containing it. A line may cite another only if the cited line lies
within its scope, that is, only if every subproof containing the cited line
also contains the citing line.
\end{definition}

The corresponding type is mutually recursive, an item being either a line or
a subproof containing further items. A rule that discharges cites a subproof,
which is written as the pair of its first and last line---the \verb|SubRef|
below---and this is exactly the pair that the Lemmon rule already names.

\begin{lstlisting}
type SubRef = (Int, Int)

data FitchItem
  = FLine Int PredFormula FitchRule
  | FSub Subproof

data Subproof = Subproof
  { subAssumeLine :: Int
  , subAssumeForm :: PredFormula
  , subBody       :: [FitchItem]
  }

data FitchRule
  = FPremise                          -- never discharged: outermost level
  | FAssume                           -- opens a subproof
  | FMP Int Int
  | FCP SubRef                        -- cites the subproof, not two lines
  | FOrE Int SubRef SubRef
  -- ... and so on
\end{lstlisting}

The same proof in Fitch:

\begin{fitchproof}
1 &                 $P$                  & Premise \\
2 & \fbar           $Q$                  & Assume \\
3 & \fbar           $P \wedge Q$         & $\wedge$I 1,2 \\
4 &                 $Q \to (P \wedge Q)$ & CP 2--3 \\
\end{fitchproof}

Two conventions deserve comment. First, an assumption never discharged is
written in Fitch as a premise at the outermost level, so the choice between
``premise'' and ``assumption'' is determined by whether the line is discharged
later; nothing in the Lemmon proof marks the difference locally. Second, we
take universal introduction to carry a side condition on the name generalised
upon, rather than to open a subproof with a flagged constant. Both conventions
are found in the literature; the first agrees with the condition already
imposed in the Lemmon system, and choosing it lets universal introduction
translate one line to one line.

\section{From Fitch to Lemmon}\label{sec:fitch-to-lemmon}

This direction is unproblematic, and it is worth seeing why, because the
reason is exactly what fails in the other direction.

\begin{definition}
Let $F$ be a Fitch proof. Define $\delta(F)$, the Lemmon proof induced by $F$,
by flattening $F$ into a sequence of lines in order and assigning to each line
$n$ the set $\Gamma(n)$ given by: $\Gamma(n) = \{n\}$ if $n$ is an assumption
or premise; $\Gamma(n) = \bigcup_{m} \Gamma(m)$ for $m$ ranging over the lines
cited, if the rule at $n$ discharges nothing; and for a discharging rule,
the same union with the discharged assumptions removed.
\end{definition}

The whole of $\delta$'s arithmetic is the following function, which computes
one line's dependency set given those of the lines before it. The
discharging cases subtract; every other rule takes the union of what it cites.

\begin{lstlisting}
depsOf :: Map Int (Set Int) -> FitchRule -> Int -> Set Int
depsOf deps r self =
  case r of
    FPremise               -> singleton self
    FAssume                -> singleton self
    FCP  (a, c)            -> delete a (look c)
    FRAA (a, c)            -> delete a (look c)
    FOrE d (a1,c1) (a2,c2) ->
      unions [ look d, delete a1 (look c1), delete a2 (look c2) ]
    FExistsE m (a, c)      -> look m `union` delete a (look c)
    _                      -> unions (map look (citedLines (toLemmonRule r)))
  where
    look n = findWithDefault empty n deps
\end{lstlisting}

\begin{theorem}\label{thm:total}
$\delta$ is total: for every Fitch proof $F$, $\delta(F)$ is a Lemmon proof,
and $F$ is correct if and only if $\delta(F)$ is.
\end{theorem}

\begin{proof}[Proof sketch]
The definition of $\Gamma$ is a recursion on the position of the line, which
terminates because a line may cite only earlier lines. That the resulting
annotation satisfies the Lemmon conditions on each rule is checked rule by
rule; the discharging rules are the only ones where anything happens, and for
those the Fitch scope condition guarantees that the discharged assumption is
available to be removed.
\end{proof}

\begin{proposition}\label{prop:minimal}
For each line $n$ of $F$, the set $\Gamma(n)$ assigned by $\delta$ is
contained in the scope of $n$, and the containment may be proper.
\end{proposition}

The containment may be proper precisely when a line sits inside a subproof
whose assumption it does not use. This has a consequence worth recording.

\begin{corollary}\label{cor:notinjective}
$\delta$ is not injective. Two Fitch proofs that differ only in the placement
of a line within subproofs it does not depend on have the same image.
\end{corollary}

Consequently the round trip from Fitch through Lemmon and back is a
normalisation rather than an identity: it returns a proof whose subproofs are
no wider than they need to be.

The next observation is elementary but, we think, has not been stated. Because
the Lemmon rules determine the dependency set of each line exactly---a system
that permitted a line to cite more assumptions than it uses would not satisfy
Proposition~\ref{prop:minimal}---the dependency column is not independent
data.

\begin{proposition}\label{prop:redundant}
Let $L$ be a Lemmon proof and let $L^{-}$ be the result of deleting its
dependency column. Then $L$ is recoverable from $L^{-}$.
\end{proposition}

\begin{proof}
The recursion in the definition of $\delta$ makes no use of the dependency
column of the source; it uses only the justifications. Applying it to $L^{-}$
returns $L$.
\end{proof}

\begin{remark}
The pedagogical function of the dependency column is thus not to record
information the reader could not otherwise obtain, but to require the student
to compute it. This is a defensible reason for a notation to carry redundant
data, and it is worth distinguishing from the reason a notation might carry
data that is genuinely needed.
\end{remark}

\section{Two obstructions in the other direction}\label{sec:obstructions}

\begin{definition}
A translation from Lemmon to Fitch is \emph{positional} if it preserves the
sequence of lines: the $n$th line of the Lemmon proof becomes the $n$th line
of the Fitch proof, with the same formula and the corresponding rule.
\end{definition}

Positional translation is what one wants, when it is available. It preserves
line numbers, so a student may read the two proofs side by side; and it
certifies that the two proofs are the same proof in two notations, rather than
two proofs of the same thing.

\begin{theorem}\label{thm:notpositional}
There are correct Lemmon proofs with no positional Fitch image.
\end{theorem}

We give two, illustrating two independent obstructions. The implementation
distinguishes them, and reports which was met and where:

\begin{lstlisting}
data TranslationError
  = NotNested    Int Int [Int]   -- line, assumption discharged, boxes open
  | OutOfScope   Int Int Int     -- line, line cited, box that closed over it
  | PremiseInBox Int Int         -- line, the box it was written inside
  -- ...
\end{lstlisting}

\begin{example}[Discharge order]\label{ex:order}
\begin{lemmonproof}
\dep{1}   & (1) & $P$                        & A \\
\dep{2}   & (2) & $Q$                        & A \\
\dep{1,2} & (3) & $P \wedge Q$               & 1,2 $\wedge$I \\
\dep{2}   & (4) & $P \to (P \wedge Q)$       & 1,3 CP \\
          & (5) & $Q \to (P \to (P \wedge Q))$ & 2,4 CP \\
\end{lemmonproof}
\end{example}

\noindent Assumption (1) is made first and discharged first. In a positional
Fitch image, the subproof opened by (1) would have to contain the subproof
opened by (2), since (1) is opened first and subproofs are properly nested;
but a subproof may be closed only when every subproof it contains has been
closed, and (4) closes (1) while (2) remains open. Fitch permits discharge
only in last-in-first-out order, and Lemmon imposes no such restriction.

\begin{example}[Positional trapping]\label{ex:trap}
\begin{lemmonproof}
\dep{1} & (1) & $P$                                     & A \\
\dep{2} & (2) & $Q$                                     & A \\
\dep{1} & (3) & $P \vee R$                              & 1 $\vee$I \\
\dep{2} & (4) & $Q \wedge Q$                            & 2,2 $\wedge$I \\
        & (5) & $Q \to (Q \wedge Q)$                    & 2,4 CP \\
\dep{1} & (6) & $(P \vee R) \wedge (Q \to (Q \wedge Q))$ & 3,5 $\wedge$I \\
\end{lemmonproof}
\end{example}

\noindent Here the discharge order is unobjectionable. The difficulty is that
line (3) is written between assumption (2) and its discharge at (5), and so in
any positional image lies inside the subproof opened by (2). It does not
depend on (2)---its dependency set is $\{1\}$---but that is of no help: when
the subproof closes at (5), line (3) passes out of scope, and line (6) cites
it. In the Lemmon proof nothing goes out of scope at all. Discharging an
assumption changes the dependency set of the new line; it does not affect the
availability of any earlier line.

Example~\ref{ex:trap} is the more instructive of the two, and it isolates the
second of our opening observations. The Lemmon proof does not say that line
(3) is ``inside'' anything. It is written where it is written, and that is a
fact about the page, not about the derivation. A positional translation is
obliged to read the page as though position meant something, and here it does
not.

\begin{remark}
Both examples can be repaired by permuting lines. Swapping (1) and (2) in
Example~\ref{ex:order} yields a proof in which the discharges nest; moving (3)
before (2) in Example~\ref{ex:trap} lifts it out of the subproof. Whether
permutation always suffices is Conjecture~\ref{conj:reorder} below.
\end{remark}

\section{Sharing, and two kinds of graph}\label{sec:trees}

Both obstructions concern the linear order of lines, and it is natural to ask
what becomes of them when the linear order is set aside. Doing so requires a
little graph-theoretic vocabulary, which we give in full, since it is not
universally familiar.

\subsection{Definitions}

\begin{definition}
A \emph{directed graph} consists of a set of \emph{nodes} together with a set
of ordered pairs of nodes, the \emph{edges}. We write $m \to n$ for an edge
from $m$ to $n$, and call $m$ a \emph{parent} of $n$. A \emph{path} is a
sequence of nodes each joined to the next by an edge; a \emph{cycle} is a path
returning to its starting node.
\end{definition}

\begin{definition}
A \emph{directed acyclic graph}, or \emph{DAG}, is a directed graph with no
cycles. A \emph{rooted tree} is a DAG with a distinguished node, the
\emph{root}, such that every node other than the root has exactly one parent
and every node is reachable from the root.
\end{definition}

The difference that matters here is the one in the parent count. In a tree,
each node hangs from exactly one place. In a DAG, a node may have several
parents---and when it does, we say the node is \emph{shared}: it does duty in
more than one context at once, without being written down more than once.

\begin{definition}
Let $G$ be a DAG with root $r$. Its \emph{unfolding} is the tree whose nodes
are the paths from $r$, with the obvious edges. Unfolding replaces each shared
node by one copy per context in which it is used. The unfolding of a tree is
itself; the unfolding of a DAG with genuine sharing is strictly larger, and
may be exponentially so.
\end{definition}

\subsection{What each notation presents}

The object that both notations record is a \emph{derivation}: a tree whose
root is the conclusion, whose internal nodes are applications of rules, and
whose children are the premises those rules require. This is Gentzen's
natural deduction in its original tree form, and neither Lemmon nor Fitch
writes it down as a tree. Both write it as a DAG, because both write each line
once and then refer back to it, which is precisely what sharing is.

\begin{definition}
The \emph{citation graph} of a Lemmon proof has the lines as nodes, with an
edge $m \to n$ whenever line $n$ cites line $m$.
\end{definition}

\begin{proposition}\label{prop:lemmondag}
The citation graph of a Lemmon proof is a DAG, and every line is available to
be cited by every later line.
\end{proposition}

\begin{proof}
Acyclicity is immediate, since a line may cite only lines with smaller
numbers. That every earlier line is available is a matter of the notation
containing no device by which a line could cease to be. Discharging an
assumption alters the dependency set recorded on a \emph{new} line; it does
nothing to any line already written.
\end{proof}

Fitch does have such a device.

\begin{definition}
The \emph{scope tree} of a Fitch proof has the subproofs as nodes, together
with a root representing the whole proof, and an edge from each subproof to
those immediately contained in it. A line \emph{belongs} to the innermost
subproof containing it. A line $m$ is \emph{visible} to a later line $n$ when
the node $m$ belongs to lies on the path from the root to the node $n$ belongs
to.
\end{definition}

\begin{proposition}\label{prop:fitchtree}
The scope tree of a Fitch proof is a rooted tree, and the citation graph of a
Fitch proof is a DAG in which every edge $m \to n$ has $m$ visible to $n$.
\end{proposition}

We can now say exactly where the two notations part company, and it is not
where the slogan ``Lemmon is a graph, Fitch is a tree'' would put it. Both
present a derivation as a DAG; both share. The difference is that Fitch
constrains its sharing and Lemmon does not. In a Lemmon proof any line may be
shared with any later line whatever. In a Fitch proof a line may be shared
only along the paths its scope tree permits, and once a subproof closes,
everything belonging to it is permanently unavailable.

\subsection{The translation, restated}

The two obstructions of Section~\ref{sec:obstructions} can now be seen for
what they are. By Proposition~\ref{prop:fitchtree}, translating a Lemmon proof
to Fitch means finding a scope tree under which every citation the source
makes is permitted. Both examples are
cases where no scope tree does: in Example~\ref{ex:order} because the
discharges demand incompatible nestings, in Example~\ref{ex:trap} because a
line is cited after the only subproof it could have belonged to has closed.

When no scope tree suffices, one may unfold. Unfolding removes sharing, and a
DAG without sharing is a tree; a tree imposes no constraint that could be
violated, since nothing is shared to be shared wrongly.

A derivation is represented directly as a tree. Where the Lemmon
justification names line numbers, the derivation names subderivations, and a
line cited twice becomes two subtrees.

\begin{lstlisting}
data Deriv = Deriv { dForm :: PredFormula, dRule :: DRule }

data DRule
  = DAssume Int                      -- discharged by an ancestor
  | DPremise Int                     -- never discharged
  | DMP Deriv Deriv
  | DCP Int PredFormula Deriv        -- discharges the named assumption
  | DOrE Deriv Int PredFormula Deriv Int PredFormula Deriv
  -- ...
\end{lstlisting}

Translation is then two functions and a choice between them. The positional
translation is attempted first, because it preserves numbering and structure;
unfolding is the fallback, and the result records which was used.

\begin{lstlisting}
data Route = Direct | ViaTree

lemmonToFitch :: Proof -> Either TranslationError (Route, FitchProof)
lemmonToFitch prf =
  case lemmonToFitchDirect prf of
    Right fp -> Right (Direct, fp)
    Left _   -> do d <- toDerivation prf
                   pure (ViaTree, derivationToFitch d)
\end{lstlisting}

\begin{definition}[The unfolding translation]\label{def:construction}
Given a correct Lemmon proof, unfold its citation graph into a derivation
rooted at the last line. Then traverse the derivation, emitting first the
premises, at the outermost level and once each however often they are used,
and thereafter, for each node, the subderivations its rule requires followed
by the line applying it, opening a subproof at each discharge.
\end{definition}

\begin{theorem}\label{thm:refuted}
The construction of Definition~\ref{def:construction} does not always yield a
correct Fitch proof.
\end{theorem}

\begin{proof}
Consider

\begin{lemmonproof}
\dep{1}   & (1) & $\forall x\, G(x)$ & A \\
\dep{2}   & (2) & $F(a)$             & A \\
\dep{1}   & (3) & $G(a)$             & 1 $\forall$E \\
\dep{1}   & (4) & $\forall y\, G(y)$ & 3 $\forall$I \\
\dep{1,2} & (5) & $F(a) \wedge \forall y\, G(y)$ & 2,4 $\wedge$I \\
\dep{1}   & (6) & $F(a) \to (F(a) \wedge \forall y\, G(y))$ & 2,5 CP \\
\end{lemmonproof}

\noindent which is correct. Line (4) generalises on $a$, and $\Gamma(4) =
\{1\}$; the sole assumption in that set has formula $\forall x\, G(x)$, in
which $a$ does not occur. The Lemmon side condition is satisfied.

The construction returns

\begin{fitchproof}
1 &                 $\forall x\, G(x)$ & Premise \\
2 & \fbar           $F(a)$             & Assume \\
3 & \fbar           $G(a)$             & $\forall$E 1 \\
4 & \fbar           $\forall y\, G(y)$ & $\forall$I 3 \\
5 & \fbar           $F(a) \wedge \forall y\, G(y)$ & $\wedge$I 2,4 \\
6 &                 $F(a) \to (F(a) \wedge \forall y\, G(y))$ & CP 2--5 \\
\end{fitchproof}

\noindent in which line 4 applies universal introduction on $a$ while $F(a)$
stands as an undischarged assumption in its scope. That is contrary to the
side condition, so the proof is not correct.
\end{proof}

It is worth being clear that this is not an artefact of the Fitch system we
chose in Section~\ref{sec:systems}. One might think to repair matters by
restating Fitch's side condition against the assumptions a line
\emph{depends on} rather than those in scope. But a Fitch line does not
record what it depends on; that is precisely the difference between the two
notations, and the reason the dependency column of a Lemmon proof, though
redundant in the sense of Proposition~\ref{prop:redundant}, is not idle.
Fitch must state its eigenvariable condition against scope, because scope is
the only thing it has. The failure is therefore a genuine one: at this rule,
Fitch's coarser bookkeeping is not merely less informative but licenses
strictly less.

Nothing is lost in the way of provability. The sequent above is Fitch-provable
by other means, and the two systems remain equivalent in the sense of proving
the same sequents. What fails is the structure-preserving translation.

There are two ways to repair it. One is to move the offending line: emit each
line at the depth of the innermost assumption in its dependency set, so that
scope and dependency coincide, and the two side conditions with them. That
would work, but it requires the subproofs so induced to nest properly, which
is Conjecture~\ref{conj:noduplication} below and is not known.

The other is to move nothing and rename instead. It is available because the
Lemmon side condition says precisely what is needed to license a substitution.

\begin{lemma}[Renaming]\label{lem:rename}
Let $D$ be a derivation of $\varphi(a)$ whose open assumptions do not contain
the name $a$, and let $b$ be a name occurring nowhere in $D$. Then $D[b/a]$ is
a derivation of $\varphi(b)$ from the same open assumptions.
\end{lemma}

This is the usual substitution lemma for natural deduction, and we do not
prove it here---though it should be proved for the present rule set. The cases
needing care are those whose own side conditions mention names: a nested
universal introduction, whose eigenvariable is untouched when it differs from
$a$ and is simply renamed when it does not; and existential elimination, whose
witness must be fresh, a condition that renaming to a globally fresh name
preserves.

The repair follows. Wherever the construction is about to emit a universal
introduction whose eigenvariable occurs in an assumption in scope, rename the
subderivation feeding it to a name occurring nowhere in the proof. Its open
assumptions are untouched, since by the Lemmon side condition $a$ does not
occur in them. Its conclusion is untouched, since abstraction removes every
occurrence of $a$. And the new name occurs in no assumption whatever, so
\emph{a fortiori} in none in scope. Note that the unfolding shares nothing, so
renaming within one subderivation cannot disturb any other.

Applied to the proof of Theorem~\ref{thm:refuted}, a single line changes:

\begin{fitchproof}
1 &                 $\forall x\, G(x)$ & Premise \\
2 & \fbar           $F(a)$             & Assume \\
3 & \fbar           $G(b)$             & $\forall$E 1 \\
4 & \fbar           $\forall y\, G(y)$ & $\forall$I 3 \\
5 & \fbar           $F(a) \wedge \forall y\, G(y)$ & $\wedge$I 2,4 \\
6 &                 $F(a) \to (F(a) \wedge \forall y\, G(y))$ & CP 2--5 \\
\end{fitchproof}

\noindent Line 3 instantiates to $b$ rather than $a$, and line 4 generalises
on a name that occurs in no assumption anywhere. Nothing else moves.

\begin{conjecture}\label{conj:total}
Let the construction of Definition~\ref{def:construction} be modified as
above, renaming the eigenvariable of a universal introduction whenever it
occurs in an assumption in scope. So modified, it yields for every correct
Lemmon proof a well-formed and correct Fitch proof with the same conclusion.
\end{conjecture}

Two things recommend this repair over the first. It is local: no line moves,
no subproof changes shape, and the translation remains positional up to the
renaming of constants. And it is independent of
Conjecture~\ref{conj:noduplication}, since it raises no question about where
lines are placed. A proof of Conjecture~\ref{conj:total} therefore requires
the four lemmas of Section~\ref{sec:open} together with
Lemma~\ref{lem:rename}, and nothing else that is currently open.

The construction has a cost, which is the cost of unfolding: a line shared $k$
times is derived $k$ times, and in the worst case the result is exponential in
the size of the source. There is also a small point of bookkeeping. If a
subproof's conclusion is a line from outside it---as in deriving $Q \to P$
from the premise $P$---then the subproof would contain nothing at all, and
Fitch requires a reiteration step to bring the line in. Lemmon has no such
rule and needs none, by Proposition~\ref{prop:lemmondag}; and the reiterated
line is a tautological consequence of the line it repeats, so nothing is added
to the system.

\section{What remains to be proved}\label{sec:open}

\subsection{Towards the totality conjecture}

A proof of Conjecture~\ref{conj:total} requires four lemmas about the
construction of Definition~\ref{def:construction}. The first three are
routine; the fourth is not, for one rule.

\begin{enumerate}
\item[(L1)] \emph{Unfolding terminates.} A line may cite only lines with
smaller numbers, so the recursion is well-founded on the line number, and the
derivation is finite.

\item[(L2)] \emph{Citations precede uses.} By induction on the derivation:
the traversal emits every subderivation a node requires before the line
applying it, so every line number a rule cites is smaller than the number of
the line citing it.

\item[(L3)] \emph{Every citation is in scope.} By induction, using two facts
about the traversal: premises are emitted at the outermost level before
anything else, and an assumption is emitted as the first line of its own
subproof, so that uses of it within that subproof are visible to it. This is
the lemma that the implementation initially failed---it emitted premises
wherever the source had written them, which put some inside subproofs---and it
is therefore the one to state most carefully.

\item[(L4)] \emph{Each rule transfers.} For each of the twenty-one rules: if
the Lemmon rule licenses $\varphi$ from the lines it cites, the corresponding
Fitch rule licenses $\varphi$ from the corresponding lines and subproofs.
Eighteen of these are immediate, the rule being the same rule with a different
citation convention. Conditional proof, reductio, disjunction elimination and
existential elimination require checking that the subproof emitted really does
run from the assumption to the conclusion the Lemmon rule named. Universal
introduction is where the unmodified construction fails
(Theorem~\ref{thm:refuted}); for the modified construction the case reduces to
Lemma~\ref{lem:rename}, which itself wants a proof for this rule set.
\end{enumerate}

\subsection{Why universal introduction, and no other rule}\label{sec:eigen}

Theorem~\ref{thm:refuted} turned on universal introduction, and it is worth
recording why that rule and no other.

Every other rule in the system is a condition on the formulas cited and the
formula concluded. Universal introduction alone imposes a condition on the
\emph{assumptions}: $\forall x\varphi(x)$ may be inferred from $\varphi(a)$
only if $a$ is arbitrary, and arbitrariness is a matter of what the line rests
on. Lemmon can say this exactly, because a Lemmon line records exactly what it
rests on. Fitch can only say it about scope. Wherever scope properly includes
the dependency set---which by Proposition~\ref{prop:minimal} it may---the two
conditions come apart, and Theorem~\ref{thm:refuted} exhibits a proof where
they do.

That the repair is a renaming rather than a rearrangement is a point in
favour of the diagnosis. The difficulty is not that the line is in the wrong
place; it is that the \emph{name} is doing double duty, standing both as the
arbitrary individual of the generalisation and as the individual the enclosing
assumption is about. Lemmon can tell those two roles apart because it knows
which assumptions the line answers to. Fitch cannot, and so a name that is
arbitrary in the one sense is not visibly arbitrary in the other. Giving the
generalisation a name of its own removes the ambiguity without touching the
structure of the proof at all.

\subsection{Two further questions}

\begin{conjecture}\label{conj:reorder}
Every correct Lemmon proof can be permuted---preserving the requirement that a
line follow the lines it cites---into one with a positional Fitch image.
\end{conjecture}

\begin{conjecture}\label{conj:noduplication}
No translation requires duplication. That is, every correct Lemmon proof has a
Fitch image in which each line of the source occurs exactly once.
\end{conjecture}

Conjecture~\ref{conj:noduplication} is suggested by the following
consideration. If line $L$ is cited by line $M$ under a non-discharging rule,
then $\Gamma(L) \subseteq \Gamma(M)$. Placing every line at the depth of the
innermost assumption in its dependency set therefore places $L$ no deeper than
$M$, and hence within $M$'s scope. What this argument does not establish is
that the resulting family of subproofs is laminar, which is what a Fitch proof
requires, and we have not been able to settle whether it must be.

Neither conjecture is now load-bearing. An earlier version of this paper
repaired Theorem~\ref{thm:refuted} by placing lines at the depth of their
dependencies, which made Conjecture~\ref{conj:total} depend on
Conjecture~\ref{conj:noduplication}; the renaming repair does not, and the two
questions below are independent curiosities about how economical a translation
can be, rather than obstacles to its existing.

Some evidence bears on both. Examples~\ref{ex:order} and~\ref{ex:trap}
translate through the construction to Fitch proofs of five and six lines
respectively---exactly the lengths of the originals. No line was duplicated in
either case. This is what one would expect if
Conjecture~\ref{conj:noduplication} holds and the unfolding is simply doing
more work than it needs to.

\section{Machine verification}\label{sec:machine}

The translations in both directions were implemented in Haskell against the
proof checker for \emph{How Logic Works}. What follows is evidence for
Conjecture~\ref{conj:total}, not a substitute for a proof of it: every item
below is a statement about particular proofs, and the conjecture is a
statement about all of them.

Every proof in a regression corpus of $78$ cases was translated, where valid,
from Lemmon to Fitch and back. Twenty-five admitted a positional translation,
and for each the round trip returned the original proof line for line,
\emph{including its dependency sets}---an instance of
Proposition~\ref{prop:redundant}, since the sets were recomputed and not
copied. The two proofs of Section~\ref{sec:obstructions} were translated
through derivation trees; the results were submitted to the proof checker and
found correct, with the same conclusions.

Sixteen further proofs, transcribed from photographs of student work, were
translated. Fourteen were well formed; all fourteen admitted positional
translations and round-tripped exactly. The remaining two contained malformed
justifications on the page---a rule cited with the wrong number of
lines---and were rejected by the parser, as they should be.

We record the empirical finding for what it is worth: neither obstruction
occurred in any student proof. They are not artefacts, as
Section~\ref{sec:obstructions} shows, but they appear to be rare in practice.

One methodological point deserves recording, because it cost us two false
negatives. The round trip is a weaker test than it appears. Since $\delta$
recomputes dependency sets from the rules rather than reading them off the
page, a Fitch proof that is \emph{malformed}---one violating the conditions in
Definition~\ref{def:fitch}---may still return the Lemmon proof it came from,
exactly. Two proofs in the corpus did precisely this, both by placing a
premise inside a subproof, where its dependency set recomputes correctly while
the subproof appears to derive the premise from its assumption. Neither was
detected by the round trip; the first was noticed by looking at the rendered
output. The remedy is to check the conditions of the notation directly, and
separately:

\begin{lstlisting}
-- Conditions the notation imposes, which a round trip cannot detect:
--   * a premise may occur only at the outermost level;
--   * a line may cite only what is in scope.
fitchWellFormed :: FitchProof -> Maybe String
\end{lstlisting}

\noindent The moral generalises beyond this case. A test that verifies an
invariant is preserved under translation cannot, by itself, establish that
either side satisfies the conditions of its own notation.

\section{What kind of equivalence is this?}\label{sec:equivalence}

It is tempting to summarise all of this by saying that Lemmon and Fitch
notation are intertranslatable, and to leave it there. The results above
suggest that the summary conceals more than it reports.

Translation from Fitch to Lemmon is total, canonical, and not injective
(Theorem~\ref{thm:total}, Corollary~\ref{cor:notinjective}). Translation from
Lemmon to Fitch is not positional (Theorem~\ref{thm:notpositional}), though we
conjecture it is total by way of the unfolding
(Conjecture~\ref{conj:total}), the unmodified construction having been refuted
by Theorem~\ref{thm:refuted} and repaired by renaming. Composing in one order returns the original
proof; composing in the other returns a normal form. In the vocabulary of
structural relations between formalisms, Lemmon proofs are a \emph{retract} of
Fitch proofs, not an isomorphic copy of them.

The reason is that the two notations record different amounts. A Lemmon proof
carries its dependency structure exactly, and carries no positional
information at all. A Fitch proof carries its dependency structure only up to
an upper bound, and carries positional information in addition---information
which is, from the point of view of the derivation, arbitrary. Neither is a
richer notation than the other. They are differently lossy.

That is a more interesting relationship than notational variance, and it is
the sort of relationship that arises whenever two formalisms are said to say
the same thing. The case is unusually clean: both formalisms are small,
completely specified, and mechanically checkable, so the question of what
survives translation can be settled rather than argued. We suggest it is worth
having such a case in view when the same question is asked of theories, where
it cannot.

\begin{thebibliography}{9}

\bibitem{halvorson2020}
Hans Halvorson.
\emph{How Logic Works: A User's Guide}.
Princeton University Press, 2020.

\bibitem{lemmon1965}
E.~J. Lemmon.
\emph{Beginning Logic}.
Nelson, 1965.

\bibitem{fitch1952}
Frederic B. Fitch.
\emph{Symbolic Logic: An Introduction}.
Ronald Press, 1952.

\bibitem{suppes1957}
Patrick Suppes.
\emph{Introduction to Logic}.
Van Nostrand, 1957.

\end{thebibliography}

\end{document}
